\documentclass[a4paper]{article}

% Basic packages
\usepackage[fontset=fandol]{ctex}
\usepackage{amsmath, amssymb}
\usepackage{graphicx}
\usepackage[margin=2.5cm]{geometry}
\usepackage[most]{tcolorbox}
\usepackage{etoolbox}
\usepackage{listings}
\usepackage{hyperref}
\usepackage{booktabs}
\usepackage{subcaption}
\usepackage{float}
\usepackage{tikz}
\IfFileExists{pgfplots.sty}{
    \usepackage{pgfplots}
    \pgfplotsset{compat=1.18}
}{}

% Highlight boxes
\newtcolorbox{knowledgebox}[1]{
    enhanced, colback=blue!5!white, colframe=blue!75!black, colbacktitle=blue!75!black,
    coltitle=white, fonttitle=\bfseries, title=#1, attach boxed title to top left={yshift=-2mm, xshift=2mm},
    boxrule=1pt, sharp corners
}

\newtcolorbox{importantbox}[1]{
    enhanced, colback=yellow!10!white, colframe=yellow!80!black, colbacktitle=yellow!80!black,
    coltitle=black, fonttitle=\bfseries, title=#1, sharp corners
}

\newtcolorbox{warningbox}[1]{
    enhanced, colback=red!5!white, colframe=red!75!black, colbacktitle=red!75!black,
    coltitle=white, fonttitle=\bfseries, title=#1, sharp corners
}

% Listings
\lstset{
    language=Python,
    basicstyle=\ttfamily\small,
    keywordstyle=\color{blue},
    stringstyle=\color{red!60!black},
    commentstyle=\color{green!60!black},
    breaklines=true,
    frame=single,
    numbers=left,
    numberstyle=\tiny\color{gray},
    captionpos=b,
    extendedchars=false
}

% Front-page metadata
\newcommand{\notetitle}{课程笔记：[在此填写课程主题]}
\newcommand{\notetheme}{[在此填写主题]}
\newcommand{\notecoresummary}{[用一句话概括视频最核心的认知]}
\newcommand{\noteauthors}{五道口纳什 \& Codex}
\newcommand{\notedate}{\today}
\newcommand{\videosourcetitle}{[在此填写原视频名称]}
\newcommand{\videochannel}{[在此填写频道或作者]}
\newcommand{\videopublishdate}{[在此填写发布日期]}
\newcommand{\videoduration}{[在此填写视频时长]}
\newcommand{\videosubtitlesource}{[在此填写字幕或转录来源]}
\newcommand{\videourl}{}
\newcommand{\videocoverpath}{}

\begin{document}

\begin{titlepage}
\centering
{\Large 课程笔记\par}
\vspace{1.2cm}
{\huge\bfseries \notetitle\par}
\vspace{0.45cm}
{\Large\bfseries \notetheme\par}
\vspace{0.3cm}
{\large \notecoresummary\par}
\vspace{0.7cm}
{\large \noteauthors\par}
\vspace{0.3cm}
{\large \notedate\par}
\vspace{1.2cm}

\ifdefempty{\videocoverpath}{
{\small 请在生成笔记时填入视频封面图的本地路径。\par}
}{
\includegraphics[width=0.82\textwidth,height=0.45\textheight,keepaspectratio]{\videocoverpath}\par
}

\vfill
\begin{tcolorbox}[width=0.9\textwidth, colback=black!2!white, colframe=black!60, sharp corners]
\textbf{视频名称}：\videosourcetitle\par
\textbf{视频作者/频道}：\videochannel\par
\textbf{发布日期}：\videopublishdate\par
\textbf{视频时长}：\videoduration\par
\textbf{字幕/转录来源}：\videosubtitlesource\par
\textbf{视频链接}：
\ifdefempty{\videourl}{
未填写
}{
\href{\videourl}{\nolinkurl{\videourl}}
}
\end{tcolorbox}
\end{titlepage}

\tableofcontents
\newpage

%% --- 正文内容开始 --- %%

% Replace this block with the generated notes.
% Fill the metadata block above, especially \notetheme, \notecoresummary,
% \videosourcetitle, \videochannel, \videopublishdate, \videosubtitlesource,
% and \videocoverpath, so the first page matches the source-aware baseline.
% Keep images outside knowledgebox/importantbox/warningbox.
% Keep video provenance in JSON artifacts. In PDF, render compact source-time
% footnotes at the bottom of the page; Markdown/HTML may keep the reader view
% clean while preserving provenance in document.json.
% End the document with a final section such as:
% \section{总结与延伸}
% covering the speaker's closing discussion and your own synthesis.

%% --- 正文内容结束 --- %%

\end{document}
